 \documentclass[11pt]{article}

% ---------- Packages ----------
\usepackage{amsmath,amssymb,amsfonts,amsthm}
\usepackage{enumitem}
\usepackage{geometry}
\usepackage{fancyhdr}
\usepackage{xcolor}    % for \textcolor
\usepackage{amsmath}   % for \operatorname
\usepackage{amssymb}   % extra math symbols
\usepackage{amsfonts}  % extra math fonts
\usepackage{tikz}

\usepackage[colorlinks=true,linkcolor=blue,citecolor=blue,urlcolor=blue]{hyperref}

% ---------- Page Geometry ----------
\geometry{margin=1in}

% ---------- Running Header / Footer ----------
\pagestyle{fancy}
\fancyhf{}
\lhead{\textbf{CS 189/289A}}
\rhead{Fall 2025}
\cfoot{\thepage}

% ---------- Helpful Macros ----------
\newcommand{\RR}{\mathbb{R}}
\newcommand{\EE}{\mathbb{E}}
\newcommand{\PP}{\mathbb{P}}
\renewcommand{\vec}[1]{\mathbf{#1}}

\begin{document}

% ---------- Title Block ----------
\begin{center}
    {\huge\bfseries Homework 2}\\[4pt]
    \textbf{Due : October 17th at 11 : 59 pm}\\
\end{center}

\vspace{0.5em}\hrule\vspace{1em}

% ---------- Deliverables ----------
\noindent\textbf{Deliverables.} Submit a single PDF of your write-up to Gradescope \emph{HW2 Written Questions}. 
\\
\\
\noindent\textbf{Start each problem on a new page.}


\vspace{0.5em}\hrule\vspace{0.75em}


%%%%%%%%%%%%%%%%%%%%%%%%%%%%%%%%%%%%%%%%%%%%%%%%%%%%%%%%%%%%%%%%%%%%%%%%%%%%%%%%
\section*{Honor Code}

\noindent\emph{Write and sign the following statement:}

\begin{quote}\small
“I certify that all solutions in this document are entirely my own and that I have
not looked at anyone else’s solution. I have given credit to all external
sources I consulted.”
\end{quote}

\newpage

%%%%%%%%%%%%%%%%%%%%%%%%%%%%%%%%%%%%%%%%%%%%%%%%%%%%%%%%%%%%%%%%%%%%%%%%%%%%%%%%
%%%%%%%%%%%%%%%%%%%%%%%%%%%%%%%%%%%%%%%%%%%%%%%%%%%%%%%%%%%%%%%%%%%%%%%%%%%%%%%%
%%%%%%%%%%%%%%%%%%%%%%%%%%%%%%%%%%%%%%%%%%%%%%%%%%%%%%%%%%%%%%%%%%%%%%%%%%%%%%%%
\section{Sum of Residuals}
Consider a dataset $\mathcal{D} = \{(x_n, y_n)\}_{n=1}^N$ with $x_n \in \mathbb{R}^{D+1}$ and $y_n \in \mathbb{R}$.  
Here we assume that $X$ already contains a column of 1s.
We model the relationship between $x$ and $y$ using linear regression with parameters $w \in \mathbb{R}^{D+1}$:  
\[
\hat{y} = x^\top w.
\]
We have computed the solution $w^\ast$ to the least squares regression problem:
\[
    w^\ast = \arg \min_{w} \frac{1}{2}\sum_{n=1}^N \left( y_n - x_n^\top w\right)^2.
\]

The residual vector is then defined as: 
\[
r = y - X w^\ast
\]

\begin{enumerate}[label=(\alph*)]
    \item (\textit{2 pts}) Use the normal equation to show that:  $X^\top r = 0$.

    \textbf{Answer.} The normal equation gives $X^\top Xw^\ast=X^\top y$.
    Substituting $r=y-Xw^\ast$, we obtain
    \[
        X^\top r=X^\top(y-Xw^\ast)
        =X^\top y-X^\top Xw^\ast=0.
    \]

    \item (\textit{2 pts}) Use the presence of a bias term (encoded in the column of 1s in $X$) to show that the sum of the residuals is 0.

    \textbf{Answer.} Since $X$ contains a column $\mathbf{1}$ of all ones,
    the corresponding component of $X^\top r=0$ is
    \[
        \mathbf{1}^\top r=\sum_{n=1}^N r_n=0.
    \]

    \item (\textit{2 pts}) Suppose we add a point $(x_j, y_j)$ to our dataset with a very large $|y_j|$ but $x_j$ is close to the mean of $x$.
    Predict qualitatively how $w$ will change and explain why least squares is sensitive to such a point.
    How can you modify the loss to make the model more robust to outliers?

    \textbf{Answer.} Assuming the point has a large residual, it pulls the
    prediction at the feature mean toward $y_j$, primarily shifting the
    intercept. If $x_j$ equals the original feature mean exactly, the slopes
    remain unchanged when the least-squares solution is unique. If it is only
    close to the mean, the slopes can also change; a sufficiently large residual
    can make this change substantial.

    Squared loss gives a large residual a disproportionately large penalty.
    The point's gradient contribution is
    $-x_j(y_j-x_j^\top w)$, which grows without bound with the residual.
    To reduce this sensitivity, use absolute-error loss or Huber loss:
    \[
        L(w)=\sum_{n=1}^N \rho_\delta(y_n-x_n^\top w),
        \qquad
        \rho_\delta(r)=
        \begin{cases}
            \frac12 r^2, & |r|\leq\delta,\\
            \delta\bigl(|r|-\frac12\delta\bigr), & |r|>\delta,
        \end{cases}
        \quad \delta>0.
    \]
    Huber loss penalizes large residuals linearly, limiting their influence.
\end{enumerate}

\newpage


\section{Weighted Linear Regression}

Consider a dataset $\mathcal{D} = \{(x_n, y_n)\}_{n=1}^N$ with $x_n \in \mathbb{R}^{D+1}$ and $y_n \in \mathbb{R}$.  Here we assume that $X$ already contains a column of 1s.
    We model the relationship between $x$ and $y$ using linear regression with parameters $w \in \mathbb{R}^{D+1}$:  
    \[
    \hat{y} = x^\top w.
    \]
    Suppose we have nonnegative weights $\{\alpha_n\}_{n=1}^N \ge 0$ for each observation and our goal is to minimize the weighted least squares objective:
    \[
              L(w) \;=\; \sum_{i=1}^n \alpha_i (y_i - x_i^\top w)^2.
    \]
    In this problem, we will derive the $w^\ast$ that minimizes this loss.

    \begin{enumerate}[label=(\alph*)]
        \item (\textit{2 pts}) Write the weighted least-squares loss function $L(w)$ in matrix form using the diagonal matrix: 
        \[
        A = \operatorname{diag}(\alpha_1,\dots,\alpha_n)\in\mathbb{R}^{N\times N}.
        \]
        Your answer, should not contain any summations.

        \textbf{Answer.} The diagonal matrix $A$ weights the squared residuals:
        \[
            \boxed{L(w)=(y-Xw)^\top A(y-Xw).}
        \]

        \item (\textit{2 pts}) Compute the gradient of the loss and express your answer in matrix form (no summations). You can do this by working with the summation form and taking partial derivatives or by using the following gradient identities:
        \begin{align*}
            \nabla_w \left(u^\top w\right) &= u \\
            \nabla_w \left( w^\top A w  \right)&= (A + A^\top) w  
        \end{align*}

        \textbf{Answer.} Since $A^\top=A$, expanding the loss gives
        \[
            L(w)=y^\top Ay-2w^\top X^\top Ay+w^\top X^\top AXw.
        \]
        The first term is constant in $w$, and $X^\top AX$ is symmetric.
        Applying the given gradient identities yields
        \[
            \boxed{\nabla_w L(w)=-2X^\top Ay+2X^\top AXw
            =2X^\top A(Xw-y).}
        \]
        
        \item (\textit{3 pts}) Set the gradient equal to zero, and solve for the closed-form solution for $w^\ast$ that minimizes the weighted least squares objective. 
        in terms of $X$, $y$, and $A$, assuming $X^\top A X$ is invertible.

        \textbf{Answer.} Setting the gradient to zero gives
        \[
            2X^\top A(Xw^\ast-y)=0
            \quad\Longrightarrow\quad X^\top AXw^\ast=X^\top Ay.
        \]
        Since $X^\top AX$ is invertible,
        \[
            \boxed{w^\ast=(X^\top AX)^{-1}X^\top Ay.}
        \]
        The Hessian $2X^\top AX$ is positive definite under this assumption,
        so this is the unique global minimizer.

        \item (\textit{3 pts}) Under what condition is $X^\top A X$ guaranteed to be invertible? Can we derive a simpler condition when $\alpha_i > 0$?

        \textbf{Answer.} For nonnegative weights, let
        $B=A^{1/2}X$. Then $X^\top AX=B^\top B$, and for every vector $v$,
        \[
            v^\top X^\top AXv=\|Bv\|_2^2.
        \]
        Thus $X^\top AX$ is positive definite, and hence invertible,
        exactly when $B$ has no nonzero null vector:
        \[
            \boxed{\operatorname{rank}(A^{1/2}X)=D+1.}
        \]
        Equivalently, the rows of $X$ with positive weights must form a
        matrix of full column rank. If every $\alpha_i>0$, then $A^{1/2}$
        is invertible and preserves rank, so the condition simplifies to
        \[
            \boxed{\operatorname{rank}(X)=D+1.}
        \]
        Positive weights alone do not suffice if the columns of $X$
        are linearly dependent.
    \end{enumerate}

\newpage



\section{Online K-Means}

Consider the sequential (online) version of the K-means algorithm. \hyperlink{https://en.wikipedia.org/wiki/Online_machine_learning}{Online algorithms} assume that you receive data points in a \textit{stream} rather than having all of your data already stored.
 
In online K-means, instead of using the full dataset to update cluster centers at each iteration, we update the nearest prototype \(\mu_k\) for each data point \(x_n\) as it is observed, using the rule:
    \[
    \mu_k \;\gets\; \mu_k + \frac{1}{N_k} \,(x_n - \mu_k),
    \]
where \(N_k\) is the number of points that have been assigned to cluster \(k\) so far. Each point is assigned to the cluster with the nearest center at the time it is observed.  

\begin{enumerate}[label=(\alph*)]
    \item (\textit{3 pts}) Explain how this sequential update differs from batch K-means in terms of memory and computation. \textit{Hint: How often is each data point used in regular K-means compared to online K-means?}

    \textbf{Answer.} Batch K-means typically stores the entire dataset and
    revisits every point in each iteration to recompute assignments and centers.
    The online version described here processes each point once upon arrival
    and immediately updates its nearest center. It only needs to retain the
    centers, cluster counts, and current point, rather than all past data.
    For $N$ points in $d$ dimensions and $K$ clusters, online storage is
    $O(Kd+K+d)$, compared with $O(Nd+Kd)$ for a typical batch implementation.
    With direct nearest-center search, online processing costs $O(NKd)$ for
    one pass, while $T$ batch iterations cost $O(TNKd)$.
    
    \item (\textit{5 pts}) Show that after each update, the objective function
    \[
    E(\{\mu_k\}) = \sum_{n=1}^N \min_k \lVert x_n - \mu_k \rVert^2
    \]
    does not increase.  

    \textbf{Answer.} As stated, this claim does not hold for the online update
    and the objective evaluated on the full dataset. Consider $K=1$ and the
    one-dimensional stream $x_1=0$, $x_2=2$, $x_3=-2$. After processing $x_1$,
    the center is $\mu=0$, giving the full-data objective
    \[
        E(0)=0^2+2^2+(-2)^2=8.
    \]
    When $x_2=2$ arrives, the cluster count becomes $N_1=2$, so
    \[
        \mu'=0+\frac12(2-0)=1.
    \]
    Evaluating on the same three points gives
    \[
        E(1)=(0-1)^2+(2-1)^2+(-2-1)^2=11>8.
    \]
    Thus an online update can increase the stated total objective.

    A weaker statement does hold: the current point's loss cannot increase.
    Writing $\eta=1/N_k\in(0,1]$, the update gives
    \[
        x_n-\mu_k'=(1-\eta)(x_n-\mu_k),
        \qquad
        \|x_n-\mu_k'\|^2=(1-\eta)^2\|x_n-\mu_k\|^2
        \leq\|x_n-\mu_k\|^2.
    \]
    Since $k$ was a nearest center before the update, the current point's
    nearest-center loss also cannot increase. However, other points may move
    farther from this center, so this inequality does not establish a decrease
    in the sum over the dataset.
\end{enumerate}    
\newpage

\section{MLE vs MAP}

In this problem, we explore the difference between the Maximum Likelihood Estimate (MLE) and the Bayesian Maximum A Posteriori (MAP) estimate.
Suppose we have a coin with unknown bias $Y$. 
If we flip the coin we observe a Bernoulli random variable:
\[
X \sim \textbf{Bernoulli}(Y) \, = \, p(X=x \mid Y) = Y^x (1-Y)^{1-x}
\]
The coin will land heads ($X=1$) with probability $Y$ and tails ($X=0$) with probability $1-Y$.
We flip the coin (independently) $N$ times to collect the dataset $\mathcal{D}=\{X_n\}_{n=1}^N$.  We observe the coin lands heads $N_1 = \sum_{n=1}^N X_n$ times and tails $N_0 = N - N_1$ times. 

\vspace{2em}

For this problem, we will also introduce a prior on the probability $Y$ that the coin lands heads.  The \hyperlink{https://en.wikipedia.org/wiki/Beta_distribution}{Beta distribution} is a common prior for probabilities.  The Beta distribution has the density:
\[
p(y\mid \alpha,\beta)=\frac{1}{B(\alpha,\beta)}y^{\alpha-1}(1-y)^{\beta-1},\quad y\in(0,1),
\]
where $\alpha>0$ and $\beta>0$ are parameters of the distribution.  The function $B(\alpha,\beta)$ is the normalization function
\[
B(\alpha, \beta) = \int_{0}^{1}y^{\alpha - 1}(1-y)^{\beta - 1} dy
\]
of the Beta distribution called the  \hyperlink{https://en.wikipedia.org/wiki/Beta_function}{Beta function}.
As a prior we will need to pick values for $\alpha$ and $\beta$.  These should be based on our beliefs about the coin. 
Here we will use $\alpha=3$ and $\beta=3$:

\begin{enumerate}[label=(\alph*)]
    \item (\textit{2 pts}) Write the likelihood $p(\mathcal{D} | Y)$ of the data given $Y$.

    \textbf{Answer.} Independence of the observations gives
    \[
        p(\mathcal D\mid Y)=\prod_{n=1}^N Y^{X_n}(1-Y)^{1-X_n}
        =\boxed{Y^{N_1}(1-Y)^{N_0}.}
    \]
 
    \item (\textit{2 pts}) Compute the derivative of the log-likelihood function with respect to $Y$ and solve for the MLE estimate for $Y$.

    \textbf{Answer.} The log-likelihood and its derivative are
    \[
        \ell(Y)=N_1\log Y+N_0\log(1-Y),\qquad
        \ell'(Y)=\frac{N_1}{Y}-\frac{N_0}{1-Y}.
    \]
    For $N_1,N_0>0$, setting the derivative to zero yields
    $N_1(1-Y)=N_0Y$, hence $\boxed{\hat Y_{\mathrm{MLE}}=N_1/N}$.
    The second derivative is $-N_1/Y^2-N_0/(1-Y)^2<0$.
    If all observations are tails or heads, the maximum over $[0,1]$
    is at $0$ or $1$, respectively, consistent with this formula.
   
    
    \item (\textit{6 pts}) Write the posterior distribution $p(Y|\mathcal{D})$.  Express this posterior in the form of a Beta distribution. What are the posterior values of $\alpha$ and $\beta$ based on the prior values of $\alpha$ and $\beta$ as well as $N_0$ and $N_1$?

    \textbf{Answer.} Bayes' rule gives
    \[
        p(Y\mid\mathcal D)\propto p(\mathcal D\mid Y)p(Y)
        \propto Y^{N_1+\alpha-1}(1-Y)^{N_0+\beta-1}.
    \]
    Thus $\boxed{Y\mid\mathcal D\sim\operatorname{Beta}(\alpha+N_1,\beta+N_0)}$,
    with normalized density
    \[
        p(Y\mid\mathcal D)=
        \frac{Y^{\alpha+N_1-1}(1-Y)^{\beta+N_0-1}}
             {B(\alpha+N_1,\beta+N_0)},\quad 0<Y<1.
    \]
    For the specified prior, the posterior parameters are $N_1+3$ and $N_0+3$.
    
    \item (\textit{4 pts}) Compute the derivative of the log-posterior and solve for the MAP Estimate.

    \textbf{Answer.} Ignoring constants independent of $Y$,
    \[
        \frac{d}{dY}\log p(Y\mid\mathcal D)
        =\frac{N_1+\alpha-1}{Y}-\frac{N_0+\beta-1}{1-Y}.
    \]
    Setting this to zero gives, when both posterior parameters exceed one,
    \[
        \hat Y_{\mathrm{MAP}}=\frac{N_1+\alpha-1}{N+\alpha+\beta-2}
        =\boxed{\frac{N_1+2}{N+4}}\quad(\alpha=\beta=3).
    \]
    Here the log-posterior has strictly negative second derivative,
    so this stationary point is its unique maximum.

    \item (\textit{2 pts}) In a 1 sentence, what is the difference between the MLE and MAP estimates? Note: simply putting the your solutions to the previous parts is not sufficient. 

    \textbf{Answer.} MLE maximizes the likelihood using only the observed data,
    whereas MAP incorporates prior beliefs and maximizes the posterior density.
    
    \item  (\textit{5 pts}) How would you increase or decrease the regularization by varying $\alpha$ and $\beta$? \textit{Hint: Regularization was covered in the \hyperlink{https://drive.google.com/file/d/1s2yjJX2L93fuT9Zc96BRJ92H-D7MMrvg/view}{bias variance trade off lectures}: more regularization results in a higher bias and lower variance. }

    \textbf{Answer.} Keep the prior symmetric and set $\alpha=\beta=a\geq1$.
    Then
    \[
        \hat Y_{\mathrm{MAP}}=
        \frac{N}{N+2(a-1)}\hat Y_{\mathrm{MLE}}
        +\frac{2(a-1)}{N+2(a-1)}\frac12.
    \]
    Increasing $a$ strengthens shrinkage toward $1/2$, reducing variance
    while generally increasing bias when the true probability differs from
    $1/2$. Decreasing $a$ toward $1$ weakens this regularization; at $a=1$
    the prior is uniform and MAP agrees with MLE. Keeping $\alpha=\beta$
    preserves the prior center while changing its strength.
    
  
\end{enumerate}


\end{document}
